% =====================================================================
%  chunks/ch03k.tex —— 原书 p.134–141（PDF 141–148）
%  第 3 章 §9 配边论作为广义同调与上同调（下）；本块结束时第 3 章正文结束
% =====================================================================
% [原文 p.134]
在 Hopf 代数的情形有 Leibniz 性质，即 $S^*$ 按自然方式成为 $S$ 上的左（或右）Milnor 模；事实上，$S^*$ 还可以借助满足
\[
\rho_{\pm,a\otimes b}(x)=R_a^*(L_{s^{\pm1}b}^*(x))
\]
的表示 $\rho_\pm$，成为更大的 Hopf 代数 $S\otimes S^t$ 上的左 Milnor 模，其中 $s$ 表示 Hopf 代数的「对极」，而 $S^t$ 表示把 $S$ 的余乘「转置」、从而对极也取反所得的代数。把 $\rho_\pm$ 的作用限制到对角 $\Delta(S)\subset S\otimes S^t$ 上，就得到一个「ad-模」。作用在 $S^*$ 上（作为 Milnor 模）的算子代数在对角上的限制，导致 Hopf 代数 $S$ 的 Drinfeld「量子双」，它仍是 Hopf 代数，其对角与 $S^*$、$S^t$ 中的相同。

Buchstaber 与 Shokurov 把代数 $A^U$ 等同于实轴上固定 $0$、且在 $0$ 处一阶导数等于 $1$ 的形式微分同胚群上的微分算子代数。由此可得（这一点由 Novikov 于 1992 年指出）：代数 $A^U$ 是 Landweber--Novikov 代数 $S$ 的算子双，相应于装备了作用 $R_a^*$ 的 Milnor 模 $S^*$；这为按 Novikov 的方案（1992）建立「微分算子环的量子类比」这一一般概念提供了基础，在该方案中曾考虑 Fourier 变换情形的这种类比。这类算子双的各种「几乎 Hopf」性质，曾由 Novikov（1992）与 Buchstaber（1994）研究。算子双构造的这一特例，1992 年由 Semenov--Tyanshanskii、Faddeev 与 Alexeev 独立研究过，他们改用「Heisenberg 双」这一名称。

% [Tutor Note: 本段把上一段结尾的「算子双」落实为具体对象 $A^U$：它既是复上配边论的稳定运算代数，又可实现为形式微分同胚群上的微分算子代数。读时不必纠缠 Hopf 代数的细节，只需记住 $A^U$ 由两个子代数 $\Lambda$ 与 $S$ 生成，而两者的换位关系由 (9.8) 控制。]
可以非常有效地用复上配边论 $\Omega_U^*(\cdot)$ 与 \emph{Adams--Novikov 谱序列}（是上文 §6、§7 所述 Adams 谱序列的类比）来计算有限复形之间映射的稳定同伦类。映射 $K\to L$ 的稳定同伦类群定义为
\[
[K,L]^s=\lim_{j\to\infty}[\Sigma^jK,\Sigma^jL].
\]
令
\[
[K,L]^s_q=[\Sigma^qK,L]^s,\qquad [K,L]^s_*=\bigoplus_q[K,L]^s_q.
\]
于是由映射 $K\to L$ 的稳定同伦类，我们得到分次阿贝尔群 $[K,L]^s_*$。特别地，对 $K=L=S^0$，我们得到球面稳定同伦群的直和：
\[
[S^0,S^0]^s_*=\bigoplus_{q\ge0}\pi_q^s=\bigoplus_{q\ge0}\pi_{N+q}(S^N),\qquad N>q+1.
\]
如何计算 $[K,L]^s_*$？设 $M$、$N$ 分别表示分次 $A^U$-模 $\Omega_U^*(K)$、$\Omega_U^*(L)$。存在一个谱序列（\emph{Adams--Novikov 谱序列}）：
% [原文 p.135]
\begin{equation}
E_m^{p,q},\quad d_m:E_m^{p,q}\to E_m^{p-m,q+m-1},\quad d_m^2=0,\qquad E_{m+1}^{*,*}=H(E_m^{*,*},d_m), \tag{9.9}
\end{equation}
其第二项为
\[
E_2^{p,q}=\Ext_{A^U}^{p,q}(N,M),
\]
相伴群为
\[
G[K,L]_n^s=\sum_{p-q=n}E_\infty^{p,q}.
\]

注意 $\Ext_{A^U}^{0,q}(N,M)$ 恰是次数为 $q$ 的、$A^U$-模同态
\[
N=\Omega_U^*(K)\to\Omega_U^*(\Sigma^qL)=\Sigma^qM
\]
的集合：
\[
\Ext_{A^U}^{0,q}(N,M)\cong\Hom_{A^U}^q(N,M).
\]
这里 $\Sigma^qM$ 作为模与 $M$ 相同，只是分次平移了 $q$。

函子 $\Ext_A^{*,*}(\,,\,)$（同调代数的基础概念）的定义如下，其中 $A$ 是环（这里 $A=A^U$）。给定任一 $A^U$-模 $N$，可以作它的\emph{自由零调分解}（保持 $N$ 上的分次）；这是一个 $A^U$-模复形 $C$，形如
\[
\cdots\to C_n\xrightarrow{d}C_{n-1}\xrightarrow{d}C_{n-2}\xrightarrow{d}\cdots\xrightarrow{d}C_1\xrightarrow{d}C_0\xrightarrow{d}N\to0,
\]
其中所有 $C_j$ 都是自由 $A^U$-模（其自由生成元可与分次相容地选取），算子 $d$ 保持分次，并且处处 $\Ker d=\Image d$。于是 $N\cong C_0/\Image d$。现在作对偶复形 $C^*=\Hom_{A^U}(C,M)$：
\[
\cdots\xleftarrow{\partial^*}C_n^*\xleftarrow{\partial^*}C_{n-1}^*\xleftarrow{\partial^*}\cdots\xleftarrow{\partial^*}C_1^*\xleftarrow{\partial^*}C_0^*\leftarrow0,
\]
其中
\[
C_n^*=\bigoplus_qC_n^{*,q}=\bigoplus_q\Hom_{A^U}^q(C_n,M)=\bigoplus_q\Hom_{A^U}^{0,q}(C_n,\Sigma^qM).
\]
复形 $C^*$ 的同调群就是对象 $\Ext_{A^U}^{*,*}$：
\[
\bigoplus_q\Ext_{A^U}^{p,q}(N,M)=\bigoplus_q\Ker\partial^*/\Image\partial^*.
\]
可以证明这些群不依赖于具有上述性质的自由 $A^U$-模分解 $C$ 的选取。

% [Tutor Note: 这里只需把握两点：其一，$\Ext$ 就是「把自由零调分解对偶化再取同调」这一标准构造；其二，(9.9) 的谱序列是 Adams 谱序列的配边类比，用 $A^U$ 代替 Steenrod 代数、用 $\Omega_U^*$ 代替 $\Z/p$-上同调。原书下一段指出：对任意广义上同调论都能造出这种谱序列，而 $K$-理论没有，因为它不是定义在稳定范畴上的。]
\emph{注}. 对任意环 $A$，对象 $\Ext_A^{*,*}$ 最早由 Eilenberg--MacLane 于 1940 年代后期引入同调代数。以下定理由他们给出：\emph{若 $A=\Z[\pi]$ 是群 $\pi$ 的整群环，且 $\Z$ 取平凡的 $A$-模结构，则对任意 $A$-模 $M$ 有同构}
\[
\Ext_A^q(\Z,M)\cong H^q(K(\pi,1);M).
\]
\hfill$\square$

% [原文 p.136]
在 $K=S^0$ 的情形，Adams--Novikov 谱序列收敛于复形 $L$ 的稳定同伦群：
\[
M=\Lambda=\Omega_U^*,\qquad N=\Omega_U^*(L),
\]
\[
E_2^{p,q}=\Ext_{A^U}^{p,q}(N,\Lambda).
\]

以 $\Ext_{A^U}^{*,*}$ 为第二项的 Adams--Novikov 谱序列 (9.9) 的存在性，体现了定义在稳定同伦型与映射稳定同伦类的\emph{稳定范畴}上的广义上同调论 $h^*(\cdot)$ 的一个相当一般的范畴性质。（注意：稳定范畴不是 Grothendieck 意义下的阿贝尔范畴，尽管映射的稳定同伦类总带有分次阿贝尔群的天然结构。又注意：$K$-理论中没有这样的谱序列，因为该理论是用非稳定对象定义的。）广义上同调论 $h^*(\cdot)$ 对于解决具体问题的用处，当然取决于该理论的具体性质：稳定运算代数 $A^h$ 能被多有效地计算；对具体空间 $X$，$A^h$-模 $h^*(X)$ 的结构能被算出多少；最后，相应于理论 $h^*(\cdot)$ 的 Adams--Novikov 谱序列中有多少个非平凡微分 $d_m$（越少越好）。在后一方面，复上配边论 $\Omega_U^*(\cdot)$ 比普通 $\Z/p$-上同调论有明显优势（尤其对素数 $p>2$），而就相应代数对象的计算而言也丝毫不差。例如对球面的稳定同伦群有
\[
E_2^{*,*}=\Ext_{A^U}^{*,*}(\Lambda,\Lambda),
\]
而群 $\Ext_{A^U}^{1,*}$ 包含 Kervaire--Milnor 对群 $J\pi_{q-1}(SO)$ 的下界估计，这在普通 Adams 谱序列中是没有的。（关于 $J$-同态的定义参见上文 §8 与第 4 章 §3。）群 $\Ext_{A^U}^{2,*}$ 给出大量新信息，更高项亦然。1970 年代若干作者对球面的 Adams--Novikov 谱序列作了深远的计算（Miller、Ravenel、Wilson、Buchstaber 等人）。\footnote{原书英译者注：关于球面 Adams--Novikov 谱序列计算结果的综述，见 D. Ravenel 的书 \emph{Complex Cobordism and Stable Homotopy of Spheres}, Academic Press, 1986。}

复上配边论中的第一 Chern 类 $\sigma_1$ 有一些值得注意的性质。上文已提到，对复线丛 $\eta$，类 $\sigma(\eta)$ 是几何上配边类。对线丛 $\eta_1,\eta_2$，如何计算类 $\sigma_1(\eta_1\otimes\eta_2)$？

记 $\sigma_1(\eta_1)=u$，$\sigma_1(\eta_2)=v$。可以证明，类 $\sigma_1(\eta_1\otimes\eta_2)$ 可以表示成 $u$、$v$ 的幂级数（Novikov，1960 年代中期）：
% [原文 p.137]
\[
\sigma_1(\eta_1\otimes\eta_2)=f(u,v)=u+v+\sum_{i,j\ge1}\lambda_{ij}u^iv^j,\qquad \lambda_{ij}\in\Omega_U^*=\Lambda,
\]
并且下列交换形式群的定义条件成立：
\begin{equation}
\begin{gathered}
f(u,f(v,w))=f(f(u,v),w),\\
f(u,v)=f(v,u).
\end{gathered}\tag{9.10}
\end{equation}

设 $\eta^{-1}$ 是使 $\eta\otimes\eta^{-1}$ 为平凡线丛的复线丛。所谓 Chern 类 $\sigma_1(\eta^{-1})=\bar u$，是 $u=\sigma_1(\eta)$ 的如下形式的幂级数：
\[
\bar u=-u+\sum_{i\ge2}\mu_iu^i,\qquad \mu_i\in\Lambda,
\]
并且
\begin{equation}
f(u,\bar u)=0. \tag{9.11}
\end{equation}
几何上配边类所得的形式群 $f(u,v)$ 可用级数
\[
g(u)=\sum_{n\ge0}\CP^n\frac{u^{n+1}}{n+1}
\]
显式算出，即（Mishchenko，1960 年代后期）：
\begin{equation}
f(u,v)=g^{-1}(g(u)+g(v)). \tag{9.12}
\end{equation}

借助形式群 $f(u,v)$，可以定义（并计算）$K$-理论中用 $U_n$ 群的表示构造的 Adams 运算在 $\Omega_U^*(\cdot)$-理论中的类比。定义
\[
\psi^k(u)=\frac1k g^{-1}(kg(u)),\qquad u\in\Omega_U^2(\CP^\infty),
\]
其中 $u$ 仍如上为几何上配边类 $\sigma_1(\eta)$。我们要求运算 $\psi^k$ 具有下列与 $K$-理论中 Adams 运算类似的性质：
\begin{equation}
\begin{gathered}
\psi^{kl}=\psi^k\psi^l,\\
\psi^k(x+y)=\psi^k(x)+\psi^k(y),\\
\psi^k(x\cdot y)=\psi^k(x)\psi^k(y).
\end{gathered}\tag{9.13}
\end{equation}
由 (9.13)，运算 $\psi^k$ 在理论 $\Omega_U^*(\cdot)\otimes\Z[\tfrac1k]$ 中是良定义的（Novikov，1960 年代后期）。

% [Tutor Note: (9.10)--(9.13) 的关键是级数 $g(u)=\sum\CP^n u^{n+1}/(n+1)$：它同时充当形式群的「对数」，使 $f$ 化为加法。$\psi^k$ 就是「$k$ 倍对数再取逆」，这正是 $K$-理论中 Adams 运算 $\psi^k$ 在形式群语言下的推广；而 $\psi^k(u)=(1/k)g^{-1}(kg(u))$ 中的系数 $1/k$，说明须在系数环中允许除以 $k$，这就是 $\Z$ 对 $k$ 作局部化的原因。]
设 $k$ 是任意有理代数。在许多场合下，对环同态 $Q:\Omega_U^*\to k$ 作分类是有用的。Hirzebruch 证明了：这样的同态 $Q$ 由级数 $Q(z)=1+q_1z+q_2z^2+\cdots$ 确定（见第 4 章 §1--§3）。把 $Q$ 在 $\CP^n$ 处的值记作 $Q(\CP^n)$，则有公式（Novikov，1960 年代后期）
\[
g_Q(u)=\sum_{n\ge0}Q(\CP^n)\frac{u^{n+1}}{n+1}=Q(g(u)),
\]
% [原文 p.138]
从而
\[
Q(z)=\frac{z}{g_Q^{-1}(z)}.
\]
由此可见，决定上述形式群的级数 $g(u)$ 在这一场合也出现了。

Chern 特征标的自然类比（\emph{Chern--Dold 特征标}，其存在性最早由 Dold 建立）可以借助级数 $g(u)$ 有效地构造出来。该特征标由 $\Lambda$-代数同态
\begin{equation}
ch_U:\Omega_U^*(K)\to H^*(K;\Lambda\otimes\Q),\qquad \Lambda=\Omega_U^*, \tag{9.14}
\end{equation}
给出，它由下述要求确定：当 $K=\CP^\infty$、$u$ 为如上生成元时，
\begin{equation}
ch_U(g(u))=t\in H^2(\CP^2;\Z), \tag{9.15}
\end{equation}
其中 $t$ 是标准生成元。（特征标 $ch_U$ 的理论及其各种应用由 Buchstaber 于 1960 年代后期引入。）
% [Tutor Note: 原书 (9.15) 把取值空间印作 $H^2(\CP^2;\Z)$；按本段设定 $K=\CP^\infty$，一般应为 $H^2(\CP^\infty;\Z)$。由于 $H^2(\CP^2;\Z)=H^2(\CP^\infty;\Z)$，就 $2$ 维上同调而言两者一致，故此处照原文保留，读者不必视为笔误之外的新定义。]

1960 年代末，Quillen 证明了：上述几何上配边类的形式群，实际上是一维形式群理论中万有 Lazard 形式群的几何实现。他得以有效地计算代数 $A^U\otimes\Z_{(p)}$（其中 $\Z_{(p)}$ 是 $\Z$ 在素数 $p$ 处的局部化）中的重要投影算子，从而完成了「约化」复上配边论（现称「Brown--Peterson」理论 $BP^*(\cdot)$）的计算，而理论 $\Omega_U^*(\cdot)\otimes\Q_p$ 分解为这些理论的直和。（理论 $BP^*(\cdot)$ 的存在性由 Brown 与 Peterson 于 1960 年代中期（非构造性地）建立。此类投影算子更早曾由 Novikov 构造，但不够有效。）这一途径对于计算稳定同伦论中出现的群 $\Ext_{A^U}^{*,*}\otimes\Z_{(p)}$（见上文）被证明是有效的。\footnote{原书英译者注：形式群论已被非常有效地应用于球面 Adams--Novikov 谱序列的计算（Miller、Morava、Ravenel、Wilson 等人）。还有别的上同调论同时推广 $K$-理论与复上配边论，即「Morava $K$-理论」$K(n)^*(\cdot)$；它们可用形式群论有效地表述，并在 Adams--Novikov 谱序列的计算以及稳定同伦范畴一般性质的研究中非常有用。（细节见 D. Ravenel, \emph{Complex Cobordism and Stable Homotopy of Spheres}, Academic Press, 1986。）}

后来（1970 年代初），Buchstaber 与 Novikov 把形式群的双值类比引入上配边论。整个 1970 年代，相应的代数理论由 Buchstaber 发展，他还把这套理论应用于辛上配边论与自对上配边论，得到若干应用。
% [原文 p.139]
从配边论观点考察代数拓扑的工作，已由若干作者（Dieck、Gottlieb、Becker 等人）详细展开；但我们在此不再深入这些发展。\footnote{双值形式群理论（由 Buchstaber 于 1970 年代后期至 1980 年代中期发展）为辛上配边环（原文此处误拼作 simplectic，应为 symplectic）的 Adams--Novikov 谱序列计算提供了强有力的工具（Buchstaber、Ivanovskii、Nadiradze、Vershinin）。}

1960 年代初 Conner 与 Floyd 注意到，配边论可用来获得关于有限群与紧 Lie 群在闭流形上光滑作用的各种结果。我们现考虑其中一些结果。例如设给定 $p$ 阶（$p$ 为素数）微分同胚 $T:M^{2n}\to M^{2n}$，$T^p=1$，其中 $M^{2n}$ 是闭 $U$-流形，且 $T$ 有孤立不动点 $p_1,\ldots,p_m$。$T$ 限制到 $M^{2n}$ 中以 $p_j$ 为中心的充分小球面上没有不动点，并决定一个奇异复配边，即由相应的透镜空间（作为轨道空间）到 Eilenberg--MacLane 空间 $K(\Z/p,1)$ 的一个映射：
\begin{equation}
\begin{gathered}
f_j:S_j^{2n-1}/\Z/p\to S_j^\infty/\Z/p=K(\Z/p,1),\\
w_j^T=(S_j^{2n-1}/\Z/p,f_j),\qquad w_j^T\in\Omega_{2n-1}^U(S_j^\infty/\Z/p).
\end{gathered}\tag{9.16}
\end{equation}
$w_j^T$ 的配边类只依赖于 $T$ 在点 $p_j$ 处的微分 $dT:\tau(M)_{p_j}\to\tau(M)_{p_j}$，即依赖于复矩阵 $A_j=dT_{p_j}$ 的特征值 $\lambda_1,\ldots,\lambda_n$：
\[
w_j^T=w_j^T(\lambda_1,\ldots,\lambda_n),\qquad \lambda_q(A_j)=\exp(2\pi ix_{jq}/p),\qquad x_{jq}\in(\Z/p)^*,
\]
其中 $x_{jq}$ 是模 $p$ 的非零剩余。从流形 $M^{2n}$ 中挖去以球面 $S_j^{2n-1}$ 为边界的一些开邻域，就得到 $\Z/p$ 无不动点作用其上的流形。特别地，这给出关系
\[
\sum_jw_j^T=0
\]
（在群 $\Omega_{2n-1}^U(S_j^\infty/\Z/p)$ 中）。

当 $p=2$ 时 $x_{jq}=1$，并且由群
\[
\Omega_{2n-1}^U(S_j^\infty/\Z/2)=\Omega_{2n-1}^U(\RP^\infty)
\]
的结构知识可以推出（见下文）
\[
w_j^T=a_n,\qquad 2^{n-1}a_n\ne0,\qquad 2^na_n=0,
\]
从而不动点的个数必须被 $2^n$ 整除（Conner 与 Floyd，1960 年代中期）。

% [Tutor Note: 这里的思想是「局部化」：把不动点附近挖掉后，剩下的流形带上自由 $\Z/p$-作用，其配边类必为零，于是各不动点的局部不变量（透镜空间上的配边类）之和为零。$w_j^T$ 只由 $T$ 在 $p_j$ 处的线性化 $A_j$ 决定，这就是为什么它们能写成剩余 $x_{jq}$ 的函数。]
当 $p>2$ 时情形更复杂，因为配边类 $w_j^T$（我们的不动点不变量）依赖于剩余 $x_{jq}\in(\Z/p)^*$，而这里的选择当然比前面多。原来，类 $w_j^T$ 可以有效地作为剩余的函数算出（Novikov、Mishchenko、Kasparov，1960 年代后期）：记 $w_j^T=\alpha(x_{j1},\ldots,x_{jn})$，则
\begin{equation}
\alpha(x_{j1},\ldots,x_{jn})=\prod_{q=1}^n\frac{u}{g^{-1}(x_{jq}g(u))}\cap\alpha(1,\ldots,1). \tag{9.17}
\end{equation}
% [原文 p.140]

透镜空间 $S^{2n-1}/\Z/p$ 的群 $\Omega_U^*(S^{2n-1}/\Z/p)$ 与 $\Omega_*^U(S^{2n-1}/\Z/p)$，可以借助形式群论如下显式描述。（我们指出：关于这些上配边群与配边群的最早部分结果由 Conner 与 Floyd 于 1960 年代初得到。）把空间 $S_j^{2n+1}/\Z/p\subset S_j^\infty/\Z/p$ 看作配边类 $\alpha(x_{j1},\ldots,x_{jn})\in\Omega_{2n+1}^U(S^\infty/\Z/p)$ 的代表。几何上配边类 $u\in\Omega_U^2(S^{2n-1}/\Z/p)$ 作为 $\Omega_U^*$-模生成环 $\Omega_U^*(S^{2n+1}/\Z/p)$，其关系为
\begin{equation}
\begin{gathered}
u^{n+2}=0,\qquad g^{-1}(pg(u))=0,\\
Du^k=\alpha(\underbrace{1,\ldots,1}_{n-k\ \text{个}})=\alpha_{n-k}\in\Omega_U^{2(n-k)+1}(S^{2n+1}/\Z/p),\\
\alpha_{n-k}\circ\alpha_{n-l}=\alpha_{n-k-l}.
\end{gathered}\tag{9.18}
\end{equation}

剩余集合 $\{x_{j1},\ldots,x_{jn}\}$（$j=1,\ldots,m$）若真的能作为 $M^{2n}$ 上某个 $\Z/p$-作用的不动点 $p_1,\ldots,p_m\in M^{2n}$ 的对应物而实现，则必须满足
\begin{equation}
\sum_{j=1}^m\alpha(x_{j1},\ldots,x_{jn})=0. \tag{9.19}
\end{equation}
把公式 (9.17) 代入，便由 (9.19) 得到关于 $mn$ 个变量 $x_{jq}$（$\not\equiv0\bmod p$）的方程组，它确定了可容许的剩余 $x_{jq}$ 的集合。

我们指出，群 $\Omega_{2n}^U\otimes\Z/p$ 中的配边类 $[M^{2n}]$ 也可以通过剩余 $x_{jq}$ 计算。

上述讨论借助同态 $\Omega_U^*\to\Omega_{SO}^*$ 可搬到定向上配边论，并有所简化。不过这里 $p=2$ 的情形是特殊的，需要不同的技巧。公式 (9.17) 可推广到变换的不动点集是 $M$ 的带非平凡法丛的子流形的情形。复合阶循环群作用于流形的情形也有人研究过（Mishchenko、Gusein--Zade、Krichever，1970 年代初）。

在圆周 $S^1$ 光滑作用于定向流形或 $U$-流形的情形，得到了特别有趣的结果。例如，Gusein--Zade（约 1970 年）完全描述了与「$S^1$ 在流形上的作用且不存在被整个群 $S^1$ 固定的点」相联系的配边群。
% [原文 p.141]
原来，这些配边群以自然方式包含了流形上\emph{任意}光滑圆周作用的各不动点的不变量（以及它们之间的关系）。关于 $S^1$ 在 $U$-流形上光滑作用的几个此类漂亮结果，是 Krichever 在 1970 年代前半期借助复上配边论与形式群论得到的。我们现叙述其中两个结果。

1. 设 $A_{(k)}$ 是由级数
\[
A_{(k)}=\frac{kue^u}{e^{ku}-1}
\]
确定的示性类（见第 4 章 §3）；则对每个 $U$-流形 $M^{2n}$，该示性类给出有理数 $A_{(k)}(M^{2n})$。若圆周 $S^1$ 光滑作用于 $U$-流形 $M^{2n}$，且 $k$ 整除 $c_1(M^{2n})$，则 $A_{(k)}(M^{2n})=0$。当 $k=2$ 时，数 $A_{(2)}(M^{2n})$ 与 $M^{2n}$ 的符号差重合。（后一结果由 Atiyah--Hirzebruch 借助 $Spin$-流形的 Atiyah--Bott 公式证明。）

2. 若 $S^1$ 作用于 $U$-流形 $M^{2n}$，则这一作用在不动点集子流形 $N_j^k\subset M^{2n}$（$j=1,\ldots,m$）的法丛上的线性化具有特征值
\[
\lambda_{1j},\ldots,\lambda_{n-k,j},\qquad \lambda_{qj}=\exp(2\pi in_{qj}),
\]
其中 $n_{qj}$（\emph{权}）是整数。设 $l_j$ 表示负权 $n_{qj}<0$ 的个数，并设 $T_y$ 是由 Hirzebruch 级数
\[
T_y(u)=\frac{(y+1)u}{1-\exp(-u(y+1))}-yu
\]
确定的、取值于 $y$ 的多项式的示性类。则下列公式成立：
\begin{equation}
\sum_{j=1}^my^{l_j}T_y(N_j^k)=T_y(M^{2n}). \tag{9.20}
\end{equation}

当 $y=1$ 时，$T_1$ 在流形上的值与它们的符号差重合，而上述结果在这一情形早先已由 Atiyah--Hirzebruch 用椭圆算子理论得到。令 $y=-1$ 则得到 Euler--Poincaré 示性数：$T_{-1}(M)=\chi(M)$，而 (9.20) 化为关于 $\chi(M)$ 的经典公式。\footnote{约 1990 年，关于圆周在流形上的作用又得到了进一步的重要结果。这里首次出现了「圈空间上的整体分析」的某些思想（Witten、Taubes）、「椭圆亏格」这一优美概念（Oshanine、Landweber），以及它们的一些推广，例如 Baker--Akhieser 函数（Krichever）。}

% CHUNK ch03k DONE
